\clearpage
\fancyhead[L]{\small CONNAISSANCE DES TEMS.}
\newcommand{\ObservedVenusFraction}{\hbox{\vrule width0pt height7.600000bp depth0pt\special{pdf:bcontent}\special{pdf:code q 0 g 0.205405 4.313514 m 0.000000 4.108108 l 0.000000 3.902703 l 0.205405 3.697297 l 0.821622 3.697297 l 1.027027 3.902703 l 1.232432 3.902703 l 1.437838 3.697297 l 1.437838 3.286486 l 1.232432 3.081081 l 1.027027 3.081081 l 0.821622 2.875676 l 0.821622 2.054054 l 1.232432 1.643243 l 1.232432 1.232432 l 1.027027 1.027027 l 1.027027 0.205405 l 1.232432 0.000000 l 1.437838 0.000000 l 1.643243 0.205405 l 1.848649 0.205405 l 2.054054 0.000000 l 2.464865 0.000000 l 2.670270 0.205405 l 2.875676 0.205405 l 3.697297 1.027027 l 3.697297 2.670270 l 3.081081 3.286486 l 3.081081 3.697297 l 3.286486 3.902703 l 3.697297 3.902703 l 3.902703 3.697297 l 4.724324 3.697297 l 4.929730 3.902703 l 4.929730 4.108108 l 4.724324 4.313514 l 3.286486 4.313514 l 3.081081 4.108108 l 1.027027 4.108108 l 0.821622 4.313514 l h 2.259459 7.600000 m 1.027027 6.367568 l 1.027027 6.162162 l 1.232432 5.956757 l 1.437838 5.956757 l 1.848649 5.545946 l 1.848649 5.135135 l 2.259459 4.724324 l 2.875676 4.724324 l 3.081081 4.929730 l 3.081081 5.135135 l 2.875676 5.340541 l 2.875676 7.394595 l 2.670270 7.600000 l h f* Q}\special{pdf:econtent}\kern4.929730bp}}

\hypertarget{native-586-lower}{}
\begin{center}\small Göttingische gelehrte Anzeigen. Stück 40. Seite 385..396. 1815 März 11.\\[3pt]\rule{.72\linewidth}{.3pt}\end{center}

\emph{Connaissance des tems, ou des mouvemens célestes, à l'usage des astronomes et des navigateurs pour
l'an} 1816; \emph{ publiée par le bureau des longitudes.} Paris. Bei der Witwe Courcier. November 1813. 360 Seiten
in gross Octav.

Wir übergehen den astronomischen Kalender und die stehenden Artikel mit Stillschweigen, um nur von den
\emph{Zusätzen} eine Anzeige zu geben, welche seit mehreren Jahren fast ganz gefehlt hatten und bei dem
vorliegenden Jahrgange wieder zahlreich und interessant sind. Zuerst ein merkwürdiger Aufsatz über die
Cometen vom Grafen \textsc{Laplace}. Dieser grosse Geometer erklärt sich hier für diejenige Hypothese,
nach welcher jene wunderbaren Weltkörper aus Verdichtung desselben Stoffs entstehen, der die Nebelsterne
bildet und im ganzen Weltraume so allgemein verbreitet zu sein scheint. Hiernach ständen also die Cometen
mit unserm Sonnensysteme nur in derselben zufälligen Beziehung, wie die
\SourceLeaf{587}Aërolithen nach der Meinung einiger Naturforscher mit der Erde. Kommen sie so weit in den Bereich unsrer
Sonne, dass die Anziehung derselben überwiegend wird, so werden sie genöthigt, elliptische oder
hyperbolische Bahnen um diese zu beschreiben. Es erklärt sich hieraus von selbst die grosse
Mannichfaltigkeit oder vielmehr die völlige Planlosigkeit in der \emph{Lage} der beobachteten
Cometenbahnen. Nur scheint eine Bedenklichkeit gegen diese Hypothese aus dem Umstande hervorzugehen, dass
wir unter allen bisher berechneten Bahnen, deren Anzahl schon weit über hundert geht, noch gar keine
entschieden hyperbolische erkannt haben, da doch die Art des Kegelschnitts, von der ursprünglichen
Geschwindigkeit des Cometen an der Stelle des Eintritts in die Sphäre der überwiegenden Sonnenanziehung
abhängig, eben so gut eine Hyperbel als eine Ellipse werden zu können scheint. Der Zweck des Aufsatzes
geht nun dahin, aus Gründen der strengen Wahrscheinlichkeitsrechnung zu beweisen, dass unter den Bahnen
aller Cometen, die uns nahe genug kommen, um beobachtet werden zu können, nur ein äusserst geringer Theil
Hyperbeln von einer so kleinen Axe sind, dass sie durch die Beobachtungen für Hyperbeln erkannt werden
können. \textsc{Laplace} findet für das Verhältniss der Wahrscheinlichkeit einer solchen Hyperbel, wo
die halbe grosse Axe 100 Halbmesser der Erdbahn nicht übersteigt, zu der Wahrscheinlichkeit der übrigen
Fälle, nemlich einer Hyperbel von grösserer Axe, einer Parabel oder einer Ellipse, den Ausdruck (I)
\[1 : \frac{\pi-2}{10}\BareRoot{\frac{r(r+200)}{2D}} - 1\]
wo $\pi$ den halben Kreisumfang für den Halbmesser 1, $D$ das Maximum des Perihelabstandes für noch zu
beobachtende Cometen, $r$ den Halbmesser der Sphäre der überwiegenden Sonnenanziehung bedeuten. Indem
er hier setzt $r = 100000$, $D = 2$, Voraussetzungen, die man füglich gelten lassen kann, findet sich, dass
man 56 gegen 1 wetten könne, dass unter hundert beobachteten Cometenbahnen keine erkennbare Hyperbel
vorkomme. Hiebei ist vorausgesetzt, dass alle Periheldistanzen von 0 bis 2 gleich möglich sind: nimmt man
an, dass für zunehmende Periheldistanzen die Wahrscheinlichkeit immer mehr abnimmt, so wird obiges
Verhältniss noch grösser, und nach einer sehr plausibeln Hypothese für die abnehmende Wahrscheinlichkeit
der zunehmenden Periheldistanzen berechnet \textsc{Laplace}, dass man 82 gegen eins auf das Nichtvorkommen
einer erkennbaren Hyperbel unter hundert beobachteten Bahnen wetten könne.

Die Wichtigkeit des Gegenstandes und der Name des Verfassers rechtfertigen das längere Verweilen bei einem
Aufsatze von nur wenigen Seiten. Das Vorstehende sollte dessen Inhalt und Zweck nur historisch berichten:
was wir noch hinzusetzen werden, ist nur für diejenigen unsrer mathematischen Leser bestimmt, die den
Aufsatz selbst studiren und vor sich haben. So gern wir der Hypothese selbst vor allen andern, die über
diese räthselhaften Himmelskörper aufgestellt sind, und namentlich vor der \textsc{Lagrange}'schen in der
Connaissance des tems 1814, den Vorzug zugestehen: so können wir doch nicht umhin, bei dem mathematischen
Theile des Aufsatzes zwei Erinnerungen zu machen. Erstlich scheint es uns eine Uebereilung zu sein, wenn
\textsc{Laplace} S. 216 sagt `la plus petite valeur de $V$ est celle qui rend nulle la quantité renfermée
sous le radical précédent.' Vielmehr ist 0 der kleinste Werth von $V$, allein die Richtigkeit des
Ausdrucks (II)
\[1{}^{\displaystyle -}\frac{\BareRoot{(1-\frac{D}{r})}}{rV}\cdot\BareRoot{(rrVV(1+\frac{D}{r})-2D)}\]
für die Wahrscheinlichkeit, dass die Periheldistanz zwischen 0 und $2D$ falle, fängt nicht von $V=0$,
sondern erst von dem von \textsc{Laplace} als kleinsten Werth bezeichneten Werthe von $V$ an; für kleinere
Werthe von $V$ ist jene Wahrscheinlichkeit nicht imaginär, wie der Ausdruck (II) geben würde, son\SourceLeaf{588}dern $= 1$, d.\,i. \emph{Gewissheit}; diese kleineren Werthe von $V$ hat \textsc{Laplace} unsrer Einsicht nach mit
Unrecht unbeachtet gelassen. Zu dem von \textsc{Laplace} entwickelten Integrale muss also noch
$1\times\int dV$ von $V = 0$ bis $V = \BareRoot{\dfrac{2D}{1+\dfrac{D}{r}}}$ genommen, d.\,i. eben
die Grösse $\BareRoot{\dfrac{2D}{1+\dfrac{D}{r}}}$ selbst, hinzugefügt werden, und hiernach würden wir,
wenn wir übrigens \textsc{Laplace}'s Entwickelungen ungeändert beibehielten, anstatt des oben gegebnen
Ausdrucks (I) folgendes Verhältniss erhalten:
\[1 : \frac{\pi}{10}\BareRoot{\frac{r(r+200)}{2D}} - 1\]
so dass man anstatt 56 wirklich 157 gegen 1 wetten könnte, was mithin für die Hypothese noch viel
günstiger wäre. Allein wir haben noch eine zweite Bemerkung zu machen, wodurch ein sehr verändertes
Resultat hervorgebracht wird. \textsc{Laplace} entwickelt S.\,217 den von ihm S.\,216 gefundenen
Integralausdruck (b) in eine Reihe, und erhält (III)
\[\frac{(\pi-2)\BareRoot{2D}}{2r} - \frac{D}{ir\BareRoot{r}}\]
wofür also nach unsrer ersten Bemerkung
\[
\frac{\pi\BareRoot{2D}}{2r} - \frac{D}{ir\BareRoot{r}}
\]
zu setzen wäre. Diese Ausdrücke nähern sich also endlichen Grenzen, indem $i$ ins unendliche wächst,
was nicht sein kann, wie man sich aus dem strengen Integral, oder selbst aus dem Differential vor
der Integration leicht überzeugen kann. In der That findet man, wenn man die Entwickelung weiter
treibt, noch hinzuzusetzen (IV)
\[+\,\frac{iDD}{2rr\BareRoot{r}}\]
Indem \textsc{Laplace} diesen Theil nicht beachtete, konnte er sein Endresultat von $U$ dadurch unabhängig
machen, dass er diese Grösse als unendlich betrachtet (welches \textsc{Laplace} eigentlich stillschweigend
dadurch voraussetzt, dass er S.\,217 $i$ unendlich gross nimmt). Allein wollte man dieselbe
Voraussetzung, dass eigentlich für die ursprünglichen Geschwindigkeiten jeder Werth ohne
Einschränkung gleich möglich sei --- noch ferner gelten lassen, indem man das Glied (IV) mit
beachtet, so würde das Endresultat nunmehr gänzlich verschieden ausfallen; es würde unendlich wenig
wahrscheinlich werden, dass unter einer endlichen Anzahl von beobachteten Cometenbahnen irgend eine
Ellipse, Parabel oder der Parabel nahe kommende Hyperbel sich befinden sollte, vielmehr würden alle
der Wahrscheinlichkeit nach von geraden Linien nicht zu unterscheidende Hyperbeln sein. Allein
offenbar ist die Voraussetzung selbst, da sie jeder endlichen Geschwindigkeit unendlich wenig
Wahrscheinlichkeit liesse, an sich unzulässig: man \emph{muss} eine endliche Grenze für $U$
annehmen; da aber nunmehr das Endresultat ganz von $U$ abhängig bleibt, und wir keinen
Entscheidungsgrund haben, darüber etwas festzusetzen, so bleibt die Aufgabe eigentlich unauflösbar.
Wenn inzwischen die Wahrscheinlichkeitsrechnung auch gleich keinen entscheidenden Beweis
\emph{für} die Hypothese liefern kann, so entscheidet sie doch, eben wegen unsrer Unwissenheit über
die Grenze $U$, auch durchaus nichts gegen die Hypothese. Jene Rechnung lehrt selbst, dass so lange
man nicht \emph{sehr} grosse Werthe für $U$ annimmt, erkennbar hyperbolische Bahnen immer sehr
selten bleiben müssen. Nach einem angestellten Ueberschlage finden wir, dass man wenigstens
$U = 1083$ annehmen müsse (die Einheit ist die mittlere Geschwindigkeit der Erde in ihrer Bahn), um
zu der Erwartung berechtigt zu sein, unter hundert Cometenbahnen eine wenigstens entschieden
hyperbolisch zu finden.

\SourceLeaf{589}Bei der Anzeige der übrigen Aufsätze werden wir uns kürzer fassen dürfen. Der gelehrte Herausgeber der C. d. T. untersucht in dem nächstfolgenden die Frage, ob \textsc{Hipparch} in Alexandrien beobachtet hat, und macht es wahrscheinlich, dass dieser grösste Astronom des Alterthums entweder niemals in Alexandrien gewesen sei, oder doch höchstens einen vorübergehenden Aufenthalt daselbst gehabt habe. --- Von demselben über die von \textsc{Tycho} zu Uranienburg gemachten Azimuthalbeobachtungen (die freilich ziemlich fehlerhaft sind, aber nur aus der ersten Zeit seines Aufenthalts auf der Insel Hveen herrühren, und uns nicht berechtigen, die richtige Aufstellung seiner Meridianinstrumente in Zweifel zu ziehen). --- Merkwürdige Auszüge aus der Handschrift von \textsc{Ptolemäus} Optik, welche sich auf der Pariser Bibliothek befindet. --- Auszug aus \textsc{Don Joseph Rodriguez} Abhandlung über die Englische Gradmessung, nebst einigen Anmerkungen. Herr \textsc{Rodriguez} suchte zu beweisen, dass wir aus den partiellen Disharmonien, welche sich bei den verschiedenen Gradmessungen gezeigt haben, doch noch nicht berechtigt sind, eine unregelmässige Gestalt des Erdkörpers als entschieden anzusehen. Hr. \textsc{Delambre} bemerkt dagegen mit Recht, dass auch das Gegentheil nicht bewiesen ist; allein was dieser Astronom in Beziehung auf \textsc{Mechain}'s Beobachtungen in Barcellona und Montjouy hinzusetzt, scheint uns durchaus das Resultat aus diesen Beobachtungen nicht zu entkräften, dass nemlich die mit den \textsc{Lenoir}'schen Kreisen bestimmten Polhöhen Ungewissheiten von mehreren Secunden übrig lassen. Um diese merkwürdigen Beobachtungen, welche zu vielem Hin- und Herreden Gelegenheit gegeben haben, gehörig würdigen zu können, muss man eigentlich die Resultate auf folgende Weise vergleichen. Der Meridianbogen zwischen Montjouy und Barcellona, welchen die geodätischen Messungen $=59''53$ gegeben haben, folgt nach den astronomischen Beobachtungen
\begin{center}
\fontsize{9}{11}\selectfont\setlength{\tabcolsep}{3pt}\begin{tabular}{@{}l c r@{}}
aus $\beta$ Zwillinge & $.\;.\;.\;.\;.\;.\;.\;.\;.\;.\;.\;.\;.\;.\;.\;.$ & $57''20$\\
$\zeta$ grosser Bär obere Culmination & $.\;.\;.\;.\;.\;.\;.\;.$ & $61.59$\\
$\zeta$ grosser Bär untere Culmination & $.\;.\;.\;.\;.\;.\;.\;.$ & $64.66$\\
$\beta$ kleiner Bär obere Culmination & $.\;.\;.\;.\;.\;.\;.\;.$ & $62.75$\\
$\beta$ kleiner Bär untere Culmination & $.\;.\;.\;.\;.\;.\;.\;.$ & $63.39$\\
$\alpha$ kleiner Bär obere Culmination & $.\;.\;.\;.\;.\;.\;.\;.$ & $62.42$\\
$\alpha$ kleiner Bär untere Culmination & $.\;.\;.\;.\;.\;.\;.\;.$ & $63.12$\\
\end{tabular}
\end{center}
Also finden wir hier unabhängig von allen Localattractionen zwischen den Resultaten aus verschiedenen Sternen Unterschiede bis zu $7''46$, die man nur der Beobachtung und dem Instrumente zuschreiben kann. --- Noch einiges über Hrn. \textsc{Olbers} bekannte Parallaxenformeln von Hrn. \textsc{Delambre}, und noch von demselben eine Anzeige seines Abrégé de l'astronomie und der Uebersetzung des Almagests von Halma. --- Verbesserte Elemente der Jupitersbewegungen von Hrn. \textsc{Bouvard}, wonach wir bald neue Tafeln zu erwarten haben. --- Verschiedenes über das Höhenmessen mit dem Barometer von Hrn. \textsc{Prony}. Nachahmung verdient eine neue hier beschriebene Art, die Genauigkeit des Ablesens zu vergrössern. Bei mehrern hier angeführten Beobachtungsreihen weichen die äussersten Resultate für den Höhenunterschied im Mittel nur 2 bis 3 Meter ab. --- Ueber die von den Spanischen Seefahrern an vielen Punkten der Erde angestellten Pendelversuche von Hrn. \textsc{Mathieu}. --- Ueber die neuen Mondstafeln von Hrn. \textsc{Burckhardt}. Was wir bei der Anzeige dieser vortrefflichen Tafeln gewünscht hatten, nemlich die vollständigen einzelnen Resultate der Vergleichung mit den Beobachtungen, finden wir hier. Merkwürdig scheint uns, dass die negativen Breitenfehler merklich überwiegend sind; wir verstehen dies so (was aber nicht ausdrücklich bemerkt ist), dass die Tafeln im Durchschnitt die Breite um $2''29$ südlicher geben, als die Beobachtungen. Dies deutet auf die Nothwendigkeit einer kleinen Correction bei den Rechnungselementen hin; wie viel davon aber etwa auf Rechnung der Constante der Parallaxe, oder des
\SourceLeaf{590}Collimationsfehlers, oder der Refraction, oder der Polhöhe komme, wird schwer zu entscheiden sein. --- Ueber die Massen der Planeten von demselben. Als Vorbereitung zu der Construction von neuen Mars- Venus- und Mercurstafeln, die wir, wie wir hieraus sehen, von diesem vortrefflichen Astronomen zu erwarten haben, hat derselbe die Elemente der Bewegung der Erde aus beinahe 4000 Beobachtungen von neuem discutirt, von welcher Arbeit wir hier die Hauptresultate erhalten. Merkwürdig ist, dass daraus eine Verminderung der Venusmasse von ungefähr $\ObservedVenusFraction$ folgt, während Herr \textsc{von Lindenau} aus den Mercursbewegungen eine ungefähr eben so starke Vergrösserung gefunden hat. --- Bemerkungen über die Vervielfältigungskreise von Hrn. \textsc{Arago}. Allen Astronomen sind die Bedenklichkeiten bekannt, welche Herr \textsc{von Zach} gegen die Kreise mit festem Niveau an der Axe aufgestellt hat. Herr \textsc{Arago} sucht dieselben in diesem Aufsatze zu beseitigen. Dass die Verbindung des Kreises mit der Säule durch die Manipulationen der geraden Beobachtungen eine kleine Veränderung erleiden könne, welche durch das Niveau nicht sichtbar werden kann, scheine zwar nach der Theorie gegründet; allein man müsse hierüber die Künstler selbst als Schiedsrichter befragen, und diese glauben nicht, dass eine merkliche Verrückung erfolgen könne, wenn die Klammer nur die gehörige Grösse, und die Alhidade auf dem Limbus eine sanfte Bewegung habe. (Nicht sowohl die Klammer, als vielmehr die Stellschraube, welche am Ende doch allein den Kreis mit der Säule verbindet, ist der Gegenstand der Bedenklichkeit. Es scheint uns, dass man über \emph{diesen} Zweifel eine Aufklärung erhalten würde, wenn man bei einigen Beobachtungsreihen eines und desselben Sterns in den geraden Beobachtungen immer wie gewöhnlich das Fernrohr durch das Zenith, bei andern Reihen hingegen mit einer entgegengesetzten Bewegung durch das Nadir auf den Stern zurückführte. Bewirkte wirklich die Bewegung des Fernrohrs eine merkliche Verrückung des Kreises selbst, so müsste man erwarten, dass die Resultate der ersten Reihen von denen der andern in einerlei Sinn abweichen.) Die von Hrn. \textsc{von Zach} zum Nachtheil der Kreise mit festem Niveau angeführten Erfahrungen seien hier desswegen nicht entscheidend, weil diese von andern Beobachtern, die so vortrefflich unter sich harmonirenden Beobachtungen an Kreisen mit beweglichem Niveau hingegen von Hrn. \textsc{von Zach} selbst herrühren. Hr. \textsc{Arago} nimmt ferner die von Hrn. \textsc{von Zach} gemissbilligte Beobachtungsmethode in Schutz, welche darin besteht, bei den geraden Beobachtungen nicht den Stand der Säule selbst zu berichtigen, sondern immer nur die Stellung der Blase gegen die Theile der Scale aufzuzeichnen und dem zu Folge das Endresultat zu corrigiren. Herr \textsc{Arago} hat den Werth der Theile der Niveauscale an dem dreifüssigen \textsc{Reichenbach}'schen Kreise, welchen die Pariser Sternwarte als Geschenk des Grafen \textsc{Laplace} besitzt, zweimal bei sehr verschiedenen Temperaturen untersucht und fast genau übereinstimmende Resultate erhalten. Inzwischen wird hiedurch doch \emph{der} Zweifel nicht gehoben, ob man berechtigt sei. das Innere der Niveauröhre für so durchaus gleichförmig gekrümmt zu halten, dass \emph{jedem einzelnen} Theile der Scale genau gleiche Werthe beizulegen sein. Wenn man überlegt, wie ganz ausserordentlich kleine Abweichungen von der regelmässigen Krümmung hier schon hinreichen, um Unterschiede von mehreren Secunden hervorzubringen, so wird man den Zweifel nicht für ungegründet halten, ob auch der allergeschickteste Künstler diese ganz vermeiden könne, und so würden wir auf alle Fälle, so lange man hierüber keine vollkommene Gewissheit hat, es für sicherer halten, jedesmal durch Correction der Säule bei den geraden Beobachtungen die Luftblase genau wieder an den Platz bei der vorhergegangenen ungeraden zu bringen.

Herr \textsc{Arago} kommt auch noch auf die merkwürdige Behauptung des Hrn. \textsc{von Zach}, dass alle Beobachtungen mit Repetitionskreisen in Rücksicht der absoluten Genauigkeit der Resultate doch immer noch eine Ungewissheit von mehreren Secunden übrig lassen, wenn sie gleich unter sich noch so schön übereinstimmen, indem \emph{verschiedne} Kreise von einander verschiedne Resultate geben. Da dieser berühmte Beobachter selbst die Quelle dieser Unterschiede noch für ein Räthsel erklärt hat, so könnte
\SourceLeaf{591}noch weniger ein anderer Astronom, der nicht einmal von allen Nebenumständen unterrichtet ist, hierüber etwas entscheiden. Herr \textsc{Arago} bringt indessen hier einige ähnliche eigene Erfahrungen bei und glaubt die Quelle davon in einem Umstande zu finden, worüber wir hier noch einiges hinzusetzen müssen. Aus 422 Beobachtungen des Polarsterns mit einem \textsc{Fortin}'schen Kreise mit festem Niveau, die aus vier unter sich vortrefflich übereinstimmenden Reihen bestanden, ergab sich die Breite von Formentera $38^\circ\ 39'\ 56''05$. Hierauf wurde das Objectiv etwas verrückt, um ein vollkommneres Sehen zu erhalten, und sieben andere wiederum vortrefflich unter sich harmonirende Reihen gaben durch 570 Beobachtungen $38^\circ\ 39'\ 53''91$. Aehnliche Unterschiede fanden sich nachher mit demselben Instrumente in Paris, wo die Polhöhe zwischen $48^\circ\ 50'\ 11''$ bis $48^\circ\ 50'\ 15''$ schwankte, je nachdem man das Objectiv so oder anders stellte, obgleich bei einer bestimmten Stellung immer äusserst genau übereinstimmende Resultate hervorgingen. Das Resultat stand immer in genauem Zusammenhange mit der Stellung des Objectivs, obwohl bei verschiedenen Beobachtern nicht auf gleiche Weise. Herr \textsc{Arago} schliesst hieraus, dass diese Anomalien lediglich von der Art den Mittelpunkt des Sterns zu schätzen herrühren. Die Fernrohre zeigen das Bild des Sterns immer viel zu gross, und Hr. \textsc{Arago} zweifelt, ob der Mittelpunkt des scheinbaren Bildes zugleich der Mittelpunkt des wahren sei. Jeder Beobachter habe nun seine eigenthümliche Weise den Mittelpunkt zu schätzen, und so erkläre sich das Phänomen auf eine sehr natürliche Art. Wir müssen gestehen, dass wir diese Erklärung nicht statthaft finden können. Wenn wir auch zugeben wollten, dass die erwähnte Verschiedenheit der beiden Mittelpunkte gegründet wäre, so müsste doch offenbar, wenn das Fernrohr halb um seine Axe gedrehet würde, derjenige Mittelpunkt, der vorher der obere war, jetzt der untere werden. Eben das muss bei der Beobachtungsart mit den Repetitionskreisen erfolgen, wo das Fernrohr bei der geraden und ungeraden Beobachtung entgegengesetzte Lagen erhält, und also die Fehler des Pointirens sich aufheben müssen. Und wenn selbst das Fernrohr ein Bild von ganz unregelmässiger Form hervorbrächte, so würde es doch ganz gleichgültig sein, auf welchen bestimmten Punkt desselben man pointirte, wenn man nur immer bei ungeraden und geraden Beobachtungen einen und denselben wählte. Darf man das Factum, so wie es berichtet wird, als ganz entschieden annehmen, so möchte man fast geneigt sein zu vermuthen, dass durch Verrückung des Objectivs die Fäden ausserhalb seines Brennpunktes kamen, also Parallaxe hatten, wo es denn freilich darauf ankam wie der Beobachter sein Auge vor das Ocular zu halten gewohnt war. Ueberhaupt ist es auffallend, dass das Fernrohr die Einrichtung hatte, das deutliche Sehen durch Verrückung des Objectivs bewirken zu müssen. Interessant sind noch in diesem Aufsatze die Resultate, welche der dreifüssige oben erwähnte \textsc{Reichenbach}'sche Kreis für die Polhöhe aus Zenithdistanzen gegeben hat, wobei man nur noch hätte wünschen können, die Rechnungselemente mitgetheilt zu sehen, damit diese Beobachtungen zugleich zur Bestimmung sowohl der Declination als Retascension des Polarsterns mit dienen könnten. Vielleicht erhalten wir indessen diese Beobachtungen ausführlicher im nächsten Bande der Connoissance des tems, für welchen auch eine Abbildung des vortrefflichen Instruments versprochen wird.
